% TN0484 / WP0185: abstract and storyline. Build from the project root:
%   ~/.claude/skills/paper-storyline/scripts/build_storyline.sh planning/tn0484_storyline.tex
% Map mode: rows follow the v0.15.0-draft22 manuscript order and numbering (TN build; the JNE build
% keeps Appendix A and moves later appendices and figures S1--S13 to the supplement). Prepared 2026-10-03.
\documentclass[11pt,a4paper,landscape]{article}
\usepackage[T1]{fontenc}
\usepackage[utf8]{inputenc}
\usepackage{lmodern,microtype}
\usepackage[left=14mm,right=14mm,top=17mm,bottom=16mm,headheight=14pt,headsep=5mm]{geometry}
\usepackage[table]{xcolor}
\usepackage{array,tabularx,xltabular,booktabs,ragged2e}
\usepackage{url}
\usepackage{fancyhdr,lastpage}
\usepackage{amsmath}
\usepackage[hidelinks]{hyperref}

% ---- paper identity ---------------------------------------------------------
\newcommand{\wpnumber}{TN0484 / WP0185}
\newcommand{\papertitle}{The cortical column as a tuned receiver}
\newcommand{\storylinedate}{21 October 2026}
\newcommand{\baselinelabel}{v0.15.0-draft22}
\newcommand{\abstractfile}{tn0484_abstract.tex}
\newcommand{\storylineauthors}{Giulio Ruffini}
% -----------------------------------------------------------------------------

\definecolor{ink}{HTML}{183C4A}
\definecolor{muted}{HTML}{53616B}
\definecolor{rulecolor}{HTML}{A6B8BF}
\definecolor{wash}{HTML}{EDF3F5}
\hypersetup{pdftitle={\wpnumber: Abstract and storyline},
  pdfauthor={\storylineauthors},
  pdfsubject={Abstract and section map of \papertitle}}
\pagestyle{fancy}
\fancyhf{}
\fancyhead[L]{\small\sffamily\color{muted}\wpnumber\quad\papertitle}
\fancyhead[R]{\small\sffamily\color{muted}Abstract and storyline\quad\storylinedate}
\fancyfoot[L]{\small\sffamily\color{muted}Internal reading map\quad\textbullet\quad Section order: \baselinelabel}
\fancyfoot[R]{\small\sffamily\color{muted}\thepage\,/\,\pageref{LastPage}}
\renewcommand{\headrulewidth}{0.3pt}
\renewcommand{\footrulewidth}{0pt}
\setlength{\parindent}{0pt}
\setlength{\parskip}{5pt}
\setlength{\tabcolsep}{3.5mm}
\renewcommand{\arraystretch}{1.12}
\arrayrulecolor{rulecolor}
\newcolumntype{L}[1]{>{\RaggedRight\arraybackslash}p{#1}}
\newcolumntype{Y}{>{\RaggedRight\arraybackslash}X}
\newcommand{\place}[2]{\textbf{\color{ink}#1}\par\smallskip #2}
% xltabular breaks across pages (tabularx does not). Its body is scanned for the literal
% \end{xltabular}, so the environment is written out in the body rather than wrapped.
\newcommand{\maphead}{%
  \rowcolor{wash}\textbf{Place in \baselinelabel} & \textbf{Scientific content to retain} &
  \textbf{Contribution to the argument} & \textbf{Examples / links}\\[3pt]
  \midrule\endhead
  \bottomrule\endlastfoot}
\newcommand{\pagetitle}[2]{%
  {\Large\sffamily\bfseries\color{ink}#1\par}
  \vspace{2pt}{\normalsize #2\par}\vspace{5pt}}

\begin{document}
\pagetitle{Three stages of temporal interference}{Draft22: prose and placement revision approved on 21 October 2026}
\IfFileExists{\abstractfile}{\input{\abstractfile}}{\input{planning/\abstractfile}}
\absObjective\ \absApproach\ \absResults\ \absSignificance

\smallskip
The approved 300-word abstract is unchanged. This pass consolidates repetition and
relocates procedures while retaining results and their local conditions. Giulio
approved the joint plan and required the full PING system to stay in Methods.
Draft20 (\texttt{e20c047}) remains the reading baseline; \texttt{24d34c1} is the
immediate editing baseline. The canonical journal and supplement are rebuilt.
The efficiency-bound block remains in its current location pending Giulio's
placement decision. No simulation or numerical result changes.
\newpage
\pagetitle{Argument and retention map}{Section order in the revised article and supplement}
\begin{tabularx}{\textwidth}{@{}L{43mm}YY@{}}
\toprule
\textbf{Place} & \textbf{Scientific content and destination} & \textbf{Contribution to the argument}\\
\midrule
\textbf{Abstract and Introduction} & A two-tone field has no difference-frequency line. Access creates local polarization; detection generates a beat and steady source; amplification determines the response. Retain the cellular and network literature, weak-perturbation scope, and approved FEM abstract sentence. & The reader learns the physical problem before the model equations. The tests distinguish the three stages rather than assigning every response to demodulation.\\[8pt]
\textbf{Theory: population curve and field coupling} & Keep the JR sigmoid, synaptic response, and direct polarization. Derive the quadratic cross term and constant term together. Define and number the dynamic coefficient; retain its weak-input, memory, and nonzero-coefficient conditions. & The rate curve first explains what mixes the carriers. The NMM1 operator notation moves to the model appendix; it is not needed to understand this step.\\[8pt]
\textbf{Theory: envelope, spatial weighting, harmonics} & Define the spatially varying envelope before using its harmonics. Retain the side-by-side map table, product/minimum relation, two-site example, fixed-pair current-ratio law, signed projections, and recurrent weighting. Keep the second-harmonic predictions, amplitude laws, and gain correction here; move the full harmonic ladder and mixed-detector inversion to supplementary Detection. & The reader can distinguish an envelope map from a detector output, then see how amplitude and harmonic measurements separate their predictions.\\[8pt]
\bottomrule\end{tabularx}
\newpage
\pagetitle{Argument and retention map}{Section order in the revised article and supplement}
\begin{tabularx}{\textwidth}{@{}L{43mm}YY@{}}
\toprule
\textbf{Place} & \textbf{Scientific content and destination} & \textbf{Contribution to the argument}\\
\midrule
\textbf{Theory: steady source} & Retain the distinction between nonnegative mean-square polarization, signed source, and circuit response. Retain the convex/concave sign, polarity invariance, beat/DC inequality, 202-fold example, the quadratic subharmonic identity and all uniformity conditions. The envelope-norm proof and orientation averages move to supplementary Spatial maps. The direct tDCS/tACS comparison becomes a separate Theory subsection with the 2--5\% result and its amplitude conventions. & The spatial restriction concerns the mean map under stated conditions; it does not prohibit an interior beat maximum. Orientation and equal-power qualifications stay here.\\[8pt]
\textbf{Theory: gain, access, operating point} & Retain the synaptic attenuation estimates, resonant susceptibility, access--detection--gain product, static versus dynamic coefficients, access-time table, cable estimate, measured versus assumed times, and PEIX identities. & After detection is understood, the reader sees what suppresses the carriers and amplifies the beat, and why this does not establish physiological access.\\[8pt]
\textbf{Theory: exact mean field} & Move the QIF introduction after the static argument. Keep the MPR equations in physical time and rate, its nonlinear equilibrium curve, threshold recruitment, and slow-synapse and carrier-relaxation conditions. Retain the signs and asymptotes of dynamic mixing. & QIF now answers a question already posed: when does a static curve suffice? A compact crossover derivation and exact coefficient belong in the article model appendix. The fixed-voltage coefficient and exponent catalogue remain in supplementary Detection. The static-rate relation and time-scale conditions stay in Theory; the repeated synaptic display moves to the article appendix.\\[8pt]
\bottomrule\end{tabularx}
\newpage
\pagetitle{Argument and retention map}{Section order in the revised article and supplement}
\begin{tabularx}{\textwidth}{@{}L{43mm}YY@{}}
\toprule
\textbf{Place} & \textbf{Scientific content and destination} & \textbf{Contribution to the argument}\\
\midrule
\textbf{Methods} & State the purpose of waveform matching, lock-in, FEM geometry, and each model comparison before settings. Put projection choices with the head model. Place the shortened circuit schematic beside the first model explanation in Theory. Keep the full PING equations, effective input, observables, and convergence summaries. Move run catalogues and extended FEM procedures to topic-organized supplementary sections. & Methods connects each test to its observable. Detailed JR run settings, PING endpoint stability, and QIF cutoff results have explicit supplementary homes.\\[8pt]
\textbf{Results: JR and FEM} & Combine the two-sentence beat demonstration with resonance. Retain the two-tone spectrum, curvature, square-law and linearization controls, kHz access checks, FEM maps and the sensitivity conclusion, pinned coupling sweep, timing ratios, and TI/tACS controls. & The evidence proceeds from beat generation to spatial consequences and recurrent amplification. A compact FEM table replaces the numerical catalogue. Sensitivity ranges and AM/two-tone comparison values move to supplementary tables; waveform identities and local limitations remain.\\[8pt]
\textbf{Results: NMM2 and spikes} & Retain dynamic-source accuracy, the two-mode coupling curve, gain exponent and finite-amplitude deviations, attracting endpoint, spike cutoff qualification, both field amplitudes, equal-RMS carrier control, sidebands, and phase statistics. & The exact model confirms mixing and gain, then the network separates beat modulation from induced gamma. Numerical diagnostics are consolidated rather than repeated in captions.\\[8pt]
\bottomrule\end{tabularx}
\newpage
\pagetitle{Argument and retention map}{Section order in the revised article and supplement}
\begin{tabularx}{\textwidth}{@{}L{43mm}YY@{}}
\toprule
\textbf{Place} & \textbf{Scientific content and destination} & \textbf{Contribution to the argument}\\
\midrule
\textbf{Discussion} & Group interpretation into access/mixing, spatial predictions, field-induced oscillations and gain, efficiency, population-model limits, and tests. Keep the Violante biological interpretation and conditions, removing repeated displacement values and relocating the ratio algebra. Move the Vieira numerical footnote to supplementary Efficiency while retaining its unmatched-input qualification. Keep the Cifra scope attached to its claim. & Each subsection answers one scientific question. The efficiency argument distinguishes generated drive from PLV and keeps the numerical ceiling within its stated family.\\[8pt]
\textbf{Predictions and Conclusion} & Retain every test row, independent amplitudes, fixed-pair steering, state dependence, harmonic calibration, timing, blockade, and equal-power controls. Put the off-resonance measurement design beside these tests. & The conclusion states the supported findings and the measurements that separate mechanisms; detailed qualifications remain beside their claims.\\[8pt]
\textbf{Notation and supplement} & Replace tilde variables with explicit physical units in the article and an explicit membrane-time convention in the uncoupled derivation. Keep all equations, citations, supplementary figures, LaNMM results, detector-bound derivation, and vocabulary metrics. & Article appendices contain notation, the JR specification, the slow-synapse display, and compact MPR crossover derivation. Supplementary sections contain access, detection, spatial maps, amplification, efficiency, laminar extension, and metrics/settings; figures accompany their topics. Article-appendix moves do not reduce the article count.\\[8pt]
\bottomrule\end{tabularx}
\newpage
\pagetitle{Changes and review conditions}{No substantive scientific removal is proposed in this pass}
The main article retains the PING system, core MPR system and static-rate curve,
access evidence table, amplitude-matching conventions, and local cutoff and
interpretation qualifications. The appendix supplies the JR specification and
compact crossover calculation. Supplementary sections hold extended procedures,
auxiliary derivations, sensitivities, and waveform checks. Each topic names the
article claim it supports. Shared source blocks prevent divergence between editions.

The approved Violante shortening preserves its mean-map ordering, biological
interpretation, and conditions. The Vieira numerical footnote moves with the
comparison details. The efficiency-bound discussion and derivation remain in their
existing narrative and supplementary locations pending Giulio's choice. Earlier
CUT? items remain retained. The approved abstract, author order, figures, simulation
caches, and numerical results are unchanged.

The first pass is recounted before further shortening is considered. Validation
covers journal and supplement builds, retained labels and citations, notation,
cross-document references, and rendered pages. Klaus checks the committed revision
against the retention list. Public GitHub, Calliope, Zenodo, and further circulation
remain held.
\end{document}
